% ============================================================
% CHAPTER 2: LITERATURE REVIEW (DETAILED SCHEMA)
% Target Length: 40–50 Pages
% Writing Principle: Intuition → Formalism → Interpretation → Computational Implication
% ============================================================

\section{Literature Review}
\label{sec:literature_review}

% ============================================================
% 2.1 COMPLETE MARKETS AND REPLICATION
% ============================================================

\subsection{Complete Markets and Replication}
\label{subsec:complete_markets}

\subsubsection{Economic Motivation and Arbitrage Pricing}

% Goals:
% - Explain arbitrage principle
% - Motivate replication as pricing mechanism
% - Introduce concept of completeness intuitively
%
% Required structure:
% 1. Economic intuition (no heavy notation)
% 2. Why replication implies unique price
% 3. Transition toward formal modeling


\subsubsection{Mathematical Market Framework}

% Goals:
% - Define filtered probability space
% - Define asset price dynamics
% - Define admissible and self-financing strategies
%
% Required structure:
% 1. Formal definitions
% 2. Minimal assumptions explanation
% 3. Short intuitive explanation of each object


\subsubsection{Fundamental Theorem of Asset Pricing}

% Goals:
% - State First FTAP (NFLVR ↔ existence of EMM)
% - State Second FTAP (uniqueness ↔ completeness)
%
% Required structure:
% 1. Formal theorem statements
% 2. Plain-language explanation
% 3. Financial interpretation
% 4. Structural limitations discussion


\subsubsection{Martingale Representation Theorem}

% Goals:
% - State MRT formally
% - Explain predictable integrand
% - Connect integrand to hedge ratio
%
% Required structure:
% 1. Formal statement
% 2. Explanation in words
% 3. Interpretation: replication as stochastic integral
% 4. Dimensional dependence discussion


\subsubsection{Replication of European and American Claims}

% Goals:
% - European payoff representation
% - Snell envelope for American claims
%
% Required structure:
% 1. Formal equation
% 2. Interpretation
% 3. Why uniqueness simplifies valuation
% 4. Computational simplicity in complete markets


\subsubsection{Structural Breakdown of Replication}

% Goals:
% - Explain why replication fails in practice
% - Introduce incompleteness
%
% Required structure:
% 1. Dimensional mismatch example
% 2. Non-traded factors example
% 3. Transition to Section 2.2


% ============================================================
% 2.2 INCOMPLETE MARKETS AND SUPER-HEDGING
% ============================================================

\subsection{Incomplete Markets and Super-Hedging}
\label{subsec:incomplete_markets}

\subsubsection{Sources of Market Incompleteness}

% Goals:
% - Structural incompleteness
% - Stochastic volatility
% - Market frictions
%
% Required structure:
% 1. Intuition first
% 2. Short formal characterization
% 3. Financial examples


\subsubsection{Non-Uniqueness of Martingale Measures}

% Goals:
% - Define set of equivalent martingale measures
% - Explain convexity
%
% Required structure:
% 1. Formal definition of M(X)
% 2. Interpretation: multiple pricing rules
% 3. Geometric intuition


\subsubsection{Super-Replication and Dual Representation}

% Goals:
% - Define super-hedging price
% - Provide dual formulation
%
% Required structure:
% 1. Formal equation
% 2. Economic interpretation
% 3. Worst-case pricing explanation
% 4. Computational difficulty


\subsubsection{American Claims in Incomplete Markets}

% Goals:
% - Nested supremum over stopping times and measures
%
% Required structure:
% 1. Formal equation
% 2. Intuitive explanation
% 3. Minimax structure discussion
% 4. Complexity implications


\subsubsection{Utility-Based and Risk-Based Alternatives}

% Goals:
% - Indifference pricing
% - Risk-adjusted pricing
%
% Required structure:
% 1. Conceptual comparison
% 2. Mathematical contrast
% 3. Relevance to thesis positioning


% ============================================================
% 2.3 OPTIONAL DECOMPOSITION OF SUPERMARTINGALES
% ============================================================

\subsection{Optional Decomposition of Supermartingales}
\label{subsec:optional_decomposition}

\subsubsection{Doob--Meyer Decomposition}

% Goals:
% - State theorem
% - Interpret martingale vs compensator
%
% Required structure:
% 1. Formal theorem
% 2. Explanation in words
% 3. Financial analogy


\subsubsection{Optional Decomposition Theorem}

% Goals:
% - State Kramkov theorem clearly
% - Explain assumptions
%
% Required structure:
% 1. Full formal statement
% 2. Interpretation of each term
% 3. Link to super-hedging


\subsubsection{Financial Interpretation}

% Goals:
% - Interpret phi as hedge
% - Interpret C as cost
%
% Required structure:
% 1. Equation
% 2. Plain explanation
% 3. Economic meaning
% 4. Structural incompleteness discussion


\subsubsection{Comparison with Martingale Representation}

% Goals:
% - Side-by-side comparison table
%
% Required structure:
% 1. Analytical contrast
% 2. Conceptual difference
% 3. Implication for computability


\subsubsection{Non-Constructive Nature and Computational Gap}

% Goals:
% - Emphasize existence vs algorithm gap
% - Motivate deep learning
%
% Required structure:
% 1. Explanation of infinite-dimensional optimization
% 2. Transition to neural approximations


% ============================================================
% 2.4 RISK MEASURES AND HEDGING CRITERIA
% ============================================================

\subsection{Risk Measures and Quantitative Hedging Criteria}
\label{subsec:risk_measures}

\subsubsection{Coherent Risk Measures}

% Required:
% - Axioms
% - Dual representation
% - Interpretation


\subsubsection{Convex Risk Measures}

% Required:
% - Generalization beyond coherence
% - Practical relevance


\subsubsection{Value-at-Risk and Conditional Value-at-Risk}

% Required:
% - Formal definitions
% - Optimization representation
% - Tail sensitivity explanation


\subsubsection{Risk Measures in Hedging Context}

% Required:
% - Shortfall risk
% - Capital interpretation
% - Regulatory perspective


% ============================================================
% 2.5 DEEP HEDGING FRAMEWORKS
% ============================================================

\subsection{Deep Hedging}
\label{subsec:deep_hedging}

\subsubsection{Backward Stochastic Differential Equations}

% Required:
% - BSDE definition
% - Y and Z interpretation
% - Control mapping


\subsubsection{Deep BSDE Solvers}

% Required:
% - Neural approximation of Z
% - Discretization idea
% - Computational structure


\subsubsection{Deep Hedging Approach}

% Required:
% - Forward simulation
% - Risk functional minimization
% - Economic interpretation


\subsubsection{Recurrent Architectures}

% Required:
% - GRU structure
% - Hidden state interpretation
% - Dynamic control analogy


\subsubsection{Feedforward Architectures}

% Required:
% - Static mapping
% - Feature-based representation
% - Limitations


\subsubsection{Limitations of Existing Literature}

% Required:
% - Focus on European claims
% - Focus on complete markets
% - Lack of optional decomposition perspective


% ============================================================
% 2.6 MARKET MODELS
% ============================================================

\subsection{Market Models}
\label{subsec:market_models}

\subsubsection{Geometric Brownian Motion}

% Required:
% - Definition
% - Completeness property
% - Hedging simplicity


\subsubsection{Multi-Driver Models}

% Required:
% - Dimensional mismatch
% - Structural incompleteness


\subsubsection{Heston Stochastic Volatility Model}

% Required:
% - Model equations
% - Non-traded volatility factor
% - Volatility risk premium


\subsubsection{Hedging Implications}

% Required:
% - Structural vs dynamic incompleteness
% - Relevance to neural methods


% ============================================================
% 2.7 PATH SIGNATURES
% ============================================================

\subsection{Path Signatures as Feature Representations}
\label{subsec:path_signatures}

\subsubsection{Definition and Algebraic Structure}

% Required:
% - Iterated integrals
% - Tensor structure
% - Intuition


\subsubsection{Universal Approximation Property}

% Required:
% - Signature as feature map
% - Theoretical importance


\subsubsection{Applications in Finance}

% Required:
% - Path-dependent derivatives
% - Volatility modeling


\subsubsection{Computational Limitations}

% Required:
% - Truncation
% - Dimensional growth


% ============================================================
% 2.8 RESEARCH GAP AND CONTRIBUTION
% ============================================================

\subsection{Research Gap and Contribution}
\label{subsec:research_gap}

% Required structure:
% 1. Synthesis of all previous sections
% 2. Clear identification of missing elements:
%    - No unified architecture comparison
%    - No optional decomposition interpretation in deep hedging
%    - No separation of structural vs dynamic incompleteness
%    - Limited CVaR-based empirical studies
% 3. Explicit statement of thesis contributions